\documentclass[10pt,a4paper]{amsart}
\usepackage[margin=19mm]{geometry}
\usepackage{amsmath,amssymb,mathtools}
\usepackage[scr=rsfs,cal=euler]{mathalfa}
\usepackage{xfrac}
\usepackage[unicode,hidelinks,bookmarks=false]{hyperref}
\newcommand{\ie}{i.e.\ }
\newcommand{\cf}{cf.\ }
\newcommand{\resp}{respectively\ }
\newcommand{\Subsection}{\subsection}
\providecommand{\nobreakdash}{\nobreak\hbox{-}}
\setlength{\emergencystretch}{2em}
\multlinegap0pt
\usepackage{amssymb}
\usepackage{mathtools}
\usepackage{tikz}
\usepackage{tcolorbox}
\newtheorem{definition}{Definition}
\newtheorem{lemma}{Lemma}
\title{Spectral duality for some modal and residuated groupoid expansions of De Morgan algebras}
\author{Joseph McDonald}
\date{Source: arXiv:2606.03389v1 (CC BY 4.0)}
\begin{document}
\maketitle

\noindent\emph{Source and licence.} Spectral duality for some modal and residuated groupoid expansions of De Morgan algebras. Version arXiv:2606.03389v1 (CC BY 4.0), distributed by arXiv under the Creative Commons Attribution 4.0 International licence, http://creativecommons.org/licenses/by/4.0/, as recorded in the arXiv licence metadata for this exact version and shown on the abstract page https://arxiv.org/abs/2606.03389v1. Full licence notice: \texttt{licenses/arxiv-2606.03389-CC-BY-4.0.txt}.\par\smallskip

\noindent\textit{Editorial scope.} Counted unit: the article's lemma spec (source0.tex lines 460--462). The article's complete original proof of it is delivered with it (source0.tex lines 463--491, one \texttt{proof} environment of 29 lines). No mathematics is written, repaired or re-derived: the statement and the proof are the article's own lines, in the article's own notation, and no retained line is substituted or re-flowed.  Retained uncounted as the source-local context the proof needs: lines 435--442.  Nothing was omitted from any retained span, so the package carries no omission record.  No retained line contains a citation, so the wrapper carries no bibliography and no bibliography entry.  No retained line contains a figure environment, an image-inclusion command or drawing code, so no image is delivered and the package declares no asset dependency.  Editorial formatting only, in addition: an AMS \texttt{amsart} preamble that restates the article's own declarations and macros used by the retained lines, namely the environments \texttt{definition}, \texttt{lemma} on separate counters, together with the standard packages those lines need.  The retained environments print in this document as Definition 1, Lemma 1 in order of appearance, whereas the article prints the same items with the numbering of its own full document.  The complete native source of the article as distributed in the arXiv e-print for this exact version is bundled unchanged as \texttt{sources/201968/source0.tex}.
\begin{definition}\label{spectral frame morphism}
    Let $X$ and $X'$ be S4 De Morgan spectral spaces. Then a function $f\colon X\to X'$ is a \emph{spectral frame morphism} provided:
    \begin{enumerate}
        \item $f$ is a spectral map; 
        \item $f(g(x))=g(f(x))$; 
        \item $f^{-1}[\nabla_RU]=\nabla_Rf^{-1}[U]$.  
    \end{enumerate}
\end{definition}

\begin{lemma}\label{homo to spec}
    Let $A$ and $A'$ be S4 De Morgan algebras and let $h\colon A\to A'$ be a homomorphism. Then the map $\mathcal{S}_1\colon\mathcal{S}_0(A')\to\mathcal{S}_0(A)$ defined by $\mathcal{S}_1(h):=h^{-1}$ is a spectral frame morphism from the S4 De Morgan spectral space $\mathcal{S}_0(A')$ to the S4 De Morgan spectral space $\mathcal{S}_0(A)$.  
\end{lemma}

\begin{proof}
To see that $\mathcal{S}_1$ is a spectral frame morphism, we first verify that $\mathcal{S}_1(h)$ is a spectral map. Hence take any $U\in\Omega(\mathcal{S}_0(A))$ where $U=\phi(a)$ for some $a\in A$. It suffices to show that $\mathcal{S}_1(h)^{-1}[\phi(a)]\in\Omega(\mathcal{S}_0(A'))$. We have: 
\begin{align*}
    \mathcal{S}_1(h)^{-1}[\phi(a)]&=(h^{-1})^{-1}[\phi(a)]\\&=\{x\in\mathfrak{P}(A):h^{-1}[x]\in\phi(a)\}
    \\&=\{x\in\mathfrak{P}(A):a\in h^{-1}[x]\}
    \\&=\{x\in\mathfrak{P}(A):h(a)\in x\}
    \\&=\phi(h(a))
\end{align*}
where $\phi(h(a))\in\Omega(\mathcal{S}_0(A'))$. To see that Definition \ref{spectral frame morphism}(2) is satisfied, observe that:  
\begin{align*}
    \mathcal{S}_1(h)[g_A(x)]&=h^{-1}[g_A(x)]\\&=\{a\in A:h(a)\in g_A(x)\}
    \\&=\{a\in A:-h(a)\not\in x\}
    \\&=\{a\in A:h(-a)\not\in x\}
    \\&=\{a\in A:-a\not\in h^{-1}[x]\}
    \\&=g_A(h^{-1}[x])
    \\&=g_A(\mathcal{S}_1(h)[x])
\end{align*}
Finally, to see that Definition \ref{spectral frame morphism}(3) note that the previously established facts that $\phi(\nabla a)=\nabla_R\phi(a)$ and $(h^{-1})^{-1}[\phi(a)]=\phi(h(a))$ together with our assumption that $h(\nabla a)=\nabla h(a)$ for any $a\in A$ gives the following: 
\begin{align*}
    \mathcal{S}_1(h)^{-1}[\nabla_R\phi(a)]&=(h^{-1})^{-1}[\nabla_R\phi(a)]
    \\&=(h^{-1})^{-1}[\phi(\nabla a)]
    \\&=\phi(h(\nabla a))
    \\&=\phi(\nabla h(a))
    \\&=\nabla_R\phi(h(a))
    \\&=\nabla_R(h^{-1})^{-1}[\phi(a)]
    \\&=\nabla_R\mathcal{S}_1(h)^{-1}[\phi(a)]
\end{align*}
This completes the proof that $\mathcal{S}_1(h)$ is a spectral frame morphism from $\mathcal{S}_0(A')$ to $\mathcal{S}_0(A)$ whenever $h\colon A\to A'$ is a homomorphism.  
\end{proof}

\end{document}
